\section{The exact adelic bridge and the original measured setting block}
\label{sec:adelic-measures-v48}

\status{This chapter supplies the entire original \texttt{MExists},
including its literal \texttt{MeasTie}, from its original antecedents.
It also supplies the generic \texttt{MGExists}, \texttt{MWCore},
\texttt{InfTransfer}, prescribed-measure extension, and unipotent
covolume-one producers.  All infinity measures supplied by the source
are retained exactly.  The arguments below construct the measures;
they do not invoke the existence of a completed global setting.
The normalized nonzero Bessel model and the Siegel--Weil orbital and
fiber comparisons remain separate original producers.}

We use the source's conjugate-linear-first Hermitian convention.
Let \(E/F\) be a CM quadratic extension, let \(\sigma\) be its
nontrivial automorphism, and let \(J\) be a nondegenerate Hermitian
matrix on \(V=E^n\).  For an \(F\)-algebra \(R\), put
\[
 E_R=R\otimes_F E,\qquad
 G(R)=\{g\in\operatorname{GL}_n(E_R):g^*Jg=J\}.
 \tag{AMD.1}
\]
Here \(g^*=\sigma(g)^t\), and the source denotes \(G(R)\) by
\texttt{HA}\((R,\sigma,J)\).  Write \(\mathbb A_F\),
\(F_\infty\), and \(\mathbb A_{F,f}\) for its full, infinite,
and finite adelic rings.  A measure on a closed subgroup is viewed
as a measure on the ambient group by the pushforward along its
inclusion; this is exactly the source operation \(\operatorname{amb}\).

\subsection{One coordinate bridge, with the rational and compact sets}

% BEGIN AMD1
\paragraph{AMD1. Tensor splitting with the actual involution and topology.}
\label{claim:AMD1}
There are canonical topological ring isomorphisms, compatible with
\(E\) and its involution,
\[
 \mathbb A_F\otimes_F E\simeq
 (F_\infty\otimes_F E)\times(\mathbb A_{F,f}\otimes_F E),
 \qquad
 F_\infty\otimes_F E\simeq\prod_{v\mid\infty}\mathbb C.
 \tag{AMD.2}
\]
For an original \(I:\texttt{AutInst}\), the second map is the
source's \(\theta_v\) at each place, and the finite map followed
by tensor interchange is its \(\tau_F\).  The induced group map
\[
 \operatorname{split}:G(\mathbb A_F)
       \longrightarrow G(F_\infty)\times G(\mathbb A_{F,f})
 \tag{AMD.3}
\]
is a topological group isomorphism and is the literal
\texttt{Measures.splitA}.

\begin{proof}
Choose \(\delta\ne0\) with \(\sigma\delta=-\delta\).
Then \(E=F\oplus F\delta\), and \(\delta^2\in F^\times\).
For every real place \(v\), \(\phi_v(\delta)\) is nonzero
and purely imaginary.  An element of the tensor algebra has the
unique expression \(a+b\delta\), with \(a,b\) in the scalar
ring.  The first map in (AMD.2) splits these two coefficients into
their infinite and finite parts.  At an infinite place it sends
\(a_v+b_v\delta\) to \(a_v+b_v\phi_v(\delta)\); the inverse
recovers \(a_v\) from the real part and \(b_v\) by division of
the imaginary part by the nonzero real number
\(\operatorname{Im}\phi_v(\delta)\).  These formulas prove
bijectivity, continuity in both directions, multiplicativity, and
compatibility with conjugation.  On a pure tensor they give
\(a\otimes x\mapsto a_v\phi_v(x)\) and
\(a\otimes x\mapsto x\otimes a_f\), precisely the two source
formulas.  The choices of \(\delta\) merely verify these canonical
maps and do not enter their definition.

Apply these maps to every matrix entry.  They preserve the equation
\(g^*Jg=J\).  Conversely, a matrix satisfying that equation has
the two-sided inverse \(J^{-1}g^*J\); thus the inverse is a
polynomial expression in its entries.  The matrix isomorphism
therefore restricts to the topological group isomorphism (AMD.3),
including for the topology that records both a matrix and its
inverse.  Its two components are exactly extension of scalars by
the two projections of \(\mathbb A_F\).  This is the definition
of \texttt{splitA}.  This argument also covers \(n=0\).
\end{proof}
% END AMD1

% BEGIN AMD2
\paragraph{AMD2. The original bridge exists and its coordinate equations determine it.}
\label{claim:AMD2}
For every original \(I:\texttt{AutInst}\) there is a unique
continuous group isomorphism
\[
 b:I.G_\infty\times I.G_f\longrightarrow
                          G(\mathbb A_F)
 \tag{AMD.4}
\]
satisfying the original \texttt{IsBridge}.  In particular the
whole original \texttt{AdBridgeIso} holds.  Its infinite and
finite factors \(b_\infty,b_f\) satisfy
\[
 \operatorname{split}\circ b=b_\infty\times b_f,
 \quad
 b_\infty(s)=\operatorname{split}(b(s,1))_\infty,
 \quad
 b_f(x)=\operatorname{split}(b(1,x))_f.
 \tag{AMD.5}
\]
They are exactly \texttt{toInfI} and \texttt{toFinI}, for every
supplied witness of \texttt{AdBridgeIso}.

\begin{proof}
The original real change-of-basis matrix satisfies
\(P_v^*\phi_v(J)P_v=J_{\epsilon_v}\).  Send a supplied real
factor \(s_v\) to \(P_vs_vP_v^{-1}\), and send its finite
factor by interchange of the two tensor factors.  AMD1 combines
these matrices into one matrix over \(\mathbb A_F\otimes_F E\).
The displayed identity for \(P_v\) shows that the real equations
\(s_v^*J_{\epsilon_v}s_v=J_{\epsilon_v}\) become exactly the
equations for \(G\).  The finite equation is unchanged by tensor
interchange.  This gives a group isomorphism.  By construction
\[
 P_v^{-1}\theta_v(b(s,x))P_v=s_v,
 \qquad \tau_F(b(s,x))=x,
 \tag{AMD.6}
\]
which are all the \texttt{IsBridge} equations.

We check the possibly unfamiliar finite topology in the original
record.  Its identity neighborhoods have the cofinal subgroups
\[
 K(m)=\{g:g\widehat L=\widehat L,
                    (g-1)\widehat L\subseteq m\widehat L\},
                    \qquad m\ge1.
 \tag{AMD.7}
\]
A finite full lattice is sandwiched, after choosing an \(E\)-basis,
between two nonzero integral scalar multiples of \(\mathcal O_E^n\).
Its completed lattice is consequently a compact open module in
the finite adelic vector space.  The sets (AMD.7) are open for the
usual restricted-product topology: at almost all places they are
the full integral lattice stabilizer, and at the finitely many
remaining places they are ordinary matrix congruence subgroups.
Conversely an arbitrary restricted-product identity neighborhood
contains (AMD.7) for an integer with sufficiently large valuations
at its finitely many exceptional rational primes.  The source
axioms \texttt{topf\_open} and \texttt{topf\_basis} say exactly
that these same subgroups form a cofinal open identity basis for
\(I.G_f\).  The two group topologies are therefore equal.  This
is the finite-topology argument of FV4, with the original lattice
and congruence subgroups kept visible.  Thus the isomorphism above
and its inverse are continuous.

Finally, all \(\theta_v\) and \(\tau_F\) together separate
the scalar algebra by AMD1.  Equations (AMD.6) therefore determine
every entry of \(b(s,x)\).  Any other source bridge is the same
map.  Splitting its coordinates proves (AMD.5), as well as the
assertion for an arbitrarily supplied witness, rather than only
the witness chosen in this proof.
\end{proof}
% END AMD2

% BEGIN AMD3
\paragraph{AMD3. Rational points and the literal maximal-compact conditions.}
\label{claim:AMD3}
The bridge satisfies the entire original \texttt{AdBridgeGF} and
\texttt{AdBridgeK}; hence the full original
\texttt{AdelicIsomorphism} holds for every \(I\).

\begin{proof}
For \(\gamma\in G(F)\), scalar extension has real coordinates
\(P_v^{-1}\phi_v(\gamma)P_v\) and finite coordinate
\(\gamma\otimes1\).  These are exactly the two conditions in
\(I.GFset\).  Conversely those conditions specify the extension
of this same \(\gamma\), by AMD1.  Thus
\(b(I.GFset)=\texttt{ratHA}\), including both inclusions and
the same rational matrix.

For the infinite compact condition, let \(s^*J_\epsilon s=J_\epsilon\).
The source requires its entries between opposite signs to vanish.
This is equivalent to \(sJ_\epsilon=J_\epsilon s\).  Together
with the unitary equation it is equivalent to \(s^*s=1\): in
the reverse direction the two preserved forms imply that \(s\)
commutes with \(J_\epsilon\).  Under \(g=PsP^{-1}\), the
identity \(s^*s=1\) is precisely
\[
 g^*(P^{-1})^*P^{-1}g=(P^{-1})^*P^{-1},
 \tag{AMD.8}
\]
which uses the source's exact majorant \(\texttt{majV}\).

For the finite condition, the completed lattice generated by
\(L\) is the product of its local completed lattices.  Here is
the finite-generation point needed for equality.  Choose finitely
many generators of \(L\).  Locally every vector in a completed
lattice is a combination of these same generators with integral
coefficients; choosing these coefficients at every place gives
integral adeles.  The resulting finite sum is exactly the original
adelic span.  The reverse inclusion is immediate by evaluation.
It follows that an adelic matrix and its inverse preserve this
span if and only if they preserve every local completed lattice.
Thus the source's \(I.Kset\) is exactly \(K_f^{HA}\) under
the tensor flip.  Combined with (AMD.8) this proves
\(b(KGA)=KHA\).  AMD2's uniqueness proves both assertions for
every source bridge, as quantified in \texttt{AdBridgeGF} and
\texttt{AdBridgeK}.  No choice of a different lattice or a
conjugate compact subgroup has entered the conclusion.
\end{proof}
% END AMD3

\subsection{The standard flag groups and their Haar products}

An original standard frame \(P=(e_1,\ldots,e_r;
f_1,\ldots,f_r;z)\) satisfies
\[
 h(e_i,e_j)=h(f_i,f_j)=0,\quad h(e_i,f_j)=\delta_{ij},
 \quad h(z,e_i)=h(z,f_i)=0,\quad h(z,z)=-T,\quad T\ne0.
 \tag{AMD.9}
\]
Put \(U=\langle e_i\rangle_E\), \(U'=\langle f_i\rangle_E\),
\(V_0=(U\oplus U')^\perp=Ez\perp W\).  These are direct
decompositions, because their indicated Gram matrices are
nondegenerate.  In particular \(2r+1\le n\).

% BEGIN AMD4
\paragraph{AMD4. The original subgroups are closed and split into their two adelic factors.}
\label{claim:AMD4}
For every frame (AMD.9), the source's groups \(N(R),H(R),S(R)\)
are closed locally compact second-countable groups for
\(R=F_v,F_\infty,\mathbb A_{F,f},\mathbb A_F\).  They are
exactly the following groups:
\[
\begin{split}
 S(R)&=\{g:ge_i-e_i\in\langle e_j:j<i\rangle_{E_R},\\
       &\hspace{37mm}gz-z\in U_R\},\\
 H(R)&=\{g:ge_i=e_i, gf_i=f_i, gz=z\},\\
 N(R)&=\{g:ge_i-e_i\in\langle e_j:j<i\rangle_{E_R},\\
       &\hspace{20mm}gv-v\in U_R\quad\hbox{for every }v\in U_R^\perp\}.
\end{split}
 \tag{AMD.10}
\]
The map (AMD.3) restricts to topological group isomorphisms
\(N(\mathbb A_F)\simeq N(F_\infty)\times N(\mathbb A_{F,f})\)
and likewise for \(H\) and \(S\).

\begin{proof}
The first two displayed conditions are the source's
\texttt{uniFlag(F0)} and \texttt{stabT(tupH)} with the ordered
tuple \(F0=(e_1,\ldots,e_r,z)\).  The third is its
\texttt{Nsub}; since \(z\in U_R^\perp\), it already implies
the extra \(z\)-condition in \(S\).  The spans and orthogonal
spaces in (AMD.10) are scalar extensions of fixed direct
summands of the finite-dimensional \(E\)-space.  Choose bases
once.  Membership then means that finitely many specified matrix
blocks vanish or are the identity.  In particular the universal
condition on \(v\) is equivalent to testing a finite basis of
\(U^\perp\).  These are closed coordinate conditions, and
they commute with the two scalar projections in (AMD.2).
The converse also holds: vanishing in both components is
vanishing in the product algebra.  This proves closedness and
all the asserted splitting isomorphisms.

The local scalar algebras are finite products of local fields;
the adelic ones are restricted products with compact open
integer modules.  They are locally compact and second countable:
there are countably many finite places, countable local bases,
and countably many choices of a finite exceptional set and local
basis elements.  The equation \(g^*Jg=J\), together with its
polynomial inverse from AMD1, makes \(G(R)\) a closed subgroup
of the corresponding matrix space.  The same properties therefore
hold for the closed groups in (AMD.10).  The argument also proves
these properties for the pointwise stabilizer of any finite
tuple, including a linearly dependent or degenerate tuple.
\end{proof}
% END AMD4

% BEGIN AMD5
\paragraph{AMD5. The exact semidirect product, without a definiteness restriction.}
\label{claim:AMD5}
For every standard frame, \(H\simeq U(W)\) and multiplication
induces an algebraic isomorphism of varieties and a homeomorphism
\[
                    N(R)\times H(R)\longrightarrow S(R),
                    \qquad(n,h)\longmapsto nh.
 \tag{AMD.11}
\]
As groups, \(S=N\rtimes H\).  This includes standard frames
for which the real group \(H\) is noncompact.

\begin{proof}
A unitary map fixing the nondegenerate space
\(U\oplus U'\oplus Ez\) pointwise preserves its orthogonal
complement \(W\), and its restriction to \(W\) is an arbitrary
unitary map.  This proves the assertion about \(H\), with the
section \(j:U(W)\to H\) acting identically on the fixed space.

An element \(s\in S\) preserves \(U\), its complete flag,
and \(U^\perp\).  It induces an isometry on
\(U^\perp/U\simeq V_0\), fixes the image of \(z\), and
therefore preserves \(W=z^\perp\) in that quotient.  Call its
restriction \(q(s)\in U(W)\).  This is a homomorphism, and
\(q\circ j\) is the identity.  By (AMD.10), its kernel is
exactly \(N\): fixing both \(z\) and \(W\) in the quotient
is the assertion that every \(v\in U^\perp\) is fixed
modulo \(U\).  Thus the inverse to (AMD.11) is
\[
       s\longmapsto\bigl(sj(q(s))^{-1},\,j(q(s))\bigr).
 \tag{AMD.12}
\]
In an adapted basis \(q(s)\) is a fixed matrix block.  Matrix
inversion on the unitary group is polynomial.  Both (AMD.11)
and (AMD.12) are consequently algebraic and continuous over
each of the scalar rings in AMD4.  All these statements follow
from \(T\ne0\); no sign condition on \(T\) or on \(W\)
has been used.
\end{proof}
% END AMD5

% BEGIN AMD6
\paragraph{AMD6. Unipotence, unimodularity, and the exact modulus of the stabilizer action.}
\label{claim:AMD6}
The groups \(N(R)\) and \(H(R)\) are unimodular.  For every
\(h\in H(R)\), conjugation by \(h\) preserves every Haar
measure on \(N(R)\).  These statements hold at finite places
and on both adelic factors, as well as at infinity.

\begin{proof}
Use the ordered blocks \(U,V_0,U'\), and reverse the order
of the dual basis when reading the last block as a complete
flag.  An element of \(N\) is upper block triangular, has
unipotent upper triangular action on \(U\), identity on
\(V_0\), and the inverse dual action on \(U'\).  It is
therefore unipotent.  Its Lie algebra consists precisely of
\[
 X=\begin{pmatrix}
 A&B&C\\0&0&-J_0^{-1}B^*\\0&0&-A^*
 \end{pmatrix},
 \qquad A\hbox{ strictly upper triangular},\qquad C^*=-C,
 \tag{AMD.13}
\]
where \(J_0\) is the Gram matrix on \(V_0\).  Indeed these
are exactly the flag equations and \(X^*J+JX=0\).
The finite polynomials \(\exp X\) and \(\log(1+Y)\) are
inverse.  They preserve the flag equations, and
\(X^*J+JX=0\) is equivalent to \((\exp X)^*J\exp X=J\).
For the reverse implication use
\(J^{-1}(\log g)^*J=\log(g^{-1})=-\log g\).
Thus these polynomials identify \(N\), as a variety, with
the nilpotent \(F\)-Lie algebra (AMD.13).

An inner automorphism of a unipotent group is unipotent on its
Lie algebra and has determinant one.  For \(H\simeq U(W)\),
the adjoint determinant is also one: after a splitting-field
extension this is conjugation by \(g\in\operatorname{GL}_m\)
on \(M_m\), whose determinant is
\((\det g)^m(\det g^{-1})^m=1\).  This polynomial identity
descends to \(F\).  Haar measure in a local analytic chart
changes under an automorphism by the absolute value of its
derivative determinant at the identity.  Applying this to inner
automorphisms proves local unimodularity of \(N\) and \(H\).

For the additional action of \(H\) on \(N\), write
\(V_0=Ez\perp W\) in (AMD.13).  Conjugation fixes the
\(A\)- and \(C\)-coordinates and the column of \(B\)
corresponding to \(z\).  On the remaining \(r\) rows of
\(B\) it acts by the dual of its action on \(W\).  Its
\(F\)-determinant is therefore
\[
       N_{E/F}(\det_E h)^{-r}=1.
 \tag{AMD.14}
\]
The last equality follows by taking determinants of
\(h^*J_Wh=J_W\).  This proves measure preservation at every
local place, including split places and residue characteristic
two.  Local Haar measures can be normalized on compact open
subgroups at almost every finite place; their restricted product
is a Haar measure on the finite adelic group.  An adelic element
belongs to these integral subgroups at almost every place, and
the local moduli just computed are one at every place.  The
restricted-product measures are therefore unimodular and the
action preserves them.  Every other Haar measure is a positive
scalar multiple, so the assertions hold for every normalization.
AMD4 gives the same conclusion on the full product.  In the
case \(r=0\), \(N=1\), and the calculation reads as the
empty determinant one.
\end{proof}
% END AMD6

% BEGIN AMD7
\paragraph{AMD7. The actual subgroup product is Haar, with no inserted scalar.}
\label{claim:AMD7}
If \(\mu_N\) and \(\mu_H\) are any given Haar measures on
\(N(R)\) and \(H(R)\), respectively, then
\[
       \mu_S=\operatorname{mult}_*(\mu_N\times\mu_H)
 \tag{AMD.15}
\]
is a Haar measure on the exact \(S(R)\).  As ambient measures,
\[
 \operatorname{amb}\mu_S
   =\operatorname{mult}_*
       (\operatorname{amb}\mu_N\times\operatorname{amb}\mu_H).
 \tag{AMD.16}
\]

\begin{proof}
Under (AMD.11), left multiplication by \(n_0h_0\) is the map
\((n,h)\mapsto(n_0h_0nh_0^{-1},h_0h)\).  AMD6 and left
invariance of the two given measures show, by product integration,
that it preserves \(\mu_N\times\mu_H\).  This proves left
invariance of (AMD.15).  It is nonzero and finite on compact
sets because (AMD.11) is a homeomorphism and a product of Haar
measures is a nonzero Radon measure finite on compact sets.
It is therefore Haar.  Composing the multiplication map with
the subgroup inclusions gives exactly multiplication in the
ambient group; functoriality of pushforwards proves (AMD.16).
There is no comparison with a separately chosen normalization
of \(\mu_S\), and hence no unrecorded multiplicative constant.
\end{proof}
% END AMD7

\subsection{Weak source measures and exact infinity transfer}

% BEGIN AMD8
\paragraph{AMD8. The source's weak Haar predicate is sufficient.}
\label{claim:AMD8}
Let \(K\) be a closed unimodular subgroup of a locally compact
second-countable group.  A nonzero, sigma-finite Borel measure
\(\nu\) on the ambient group which is supported on \(K\)
and invariant under left multiplication by every element of
\(K\) is the ambient pushforward of a genuine Haar measure
on \(K\).  These are exactly the four assumptions in the
original \texttt{IsHaarOnSet}; local finiteness need not be
added to them.

\begin{proof}
A closed embedding identifies Borel subsets of \(K\) with
ambient Borel subsets contained in \(K\).  Pull back \(\nu\)
by this identification.  Support on \(K\) proves that its
pushforward is \(\nu\); sigma-finiteness, nonzero mass, and
left invariance pass to this measure \(\nu_K\).

Choose any ordinary Haar measure \(\lambda\) on \(K\).
Second countability implies sigma-compactness, so \(\lambda\)
is sigma-finite.  Unimodularity makes it right invariant and
inversion invariant.  For any Borel sets \(A,B\subseteq K\),
Tonelli's theorem applied to the nonnegative indicator function
gives
\[
\begin{split}
 \int_K\nu_K(gA\cap B)\,d\lambda(g)
        &=\lambda(A)\nu_K(B),\\
 \int_K\nu_K(A\cap g^{-1}B)\,d\lambda(g)
        &=\lambda(B)\nu_K(A).
\end{split}
 \tag{AMD.17}
\]
For the first equality, for fixed \(x\in B\) the admissible
\(g\)'s form \(xA^{-1}\).  For the second, for fixed
\(x\in A\) they form \(Bx^{-1}\).  The two integrands on
the left are equal by left invariance of \(\nu_K\).  Thus
\[
                \lambda(A)\nu_K(B)
                     =\lambda(B)\nu_K(A).
 \tag{AMD.18}
\]
All statements so far are valid as equalities of extended
nonnegative integrals; no local finiteness of \(\nu_K\)
has been used.

Sigma-finiteness and nonzero mass supply a Borel \(A\) with
\(0<\nu_K(A)<\infty\).  Choose a compact \(B_0\) with
\(0<\lambda(B_0)<\infty\).  The right side of (AMD.18) for
these sets is finite and positive.  Consequently both
\(\lambda(A)\) and \(\nu_K(B_0)\) are finite and positive:
a zero factor would give zero, while an infinite factor paired
with a positive factor would give infinity.  Applying (AMD.18)
again, now to arbitrary \(B\), proves
\[
             \nu_K(B)=
                \frac{\nu_K(A)}{\lambda(A)}\lambda(B).
 \tag{AMD.19}
\]
This is a finite positive scalar multiple of Haar measure, and
is therefore Haar.  The proof establishes precisely the missing
regularity instead of assuming it.
\end{proof}
% END AMD8

% BEGIN AMD9
\paragraph{AMD9. The full original infinity-transfer theorem.}
\label{claim:AMD9}
The literal \texttt{InfTransfer} holds for every original
\(I\), every supplied bridge witness, every standard frame,
and every pair of measures satisfying its two original
\texttt{IsHaarOnSet} predicates.  Its conclusion retains the
given ambient measures exactly.

\begin{proof}
Under \(b_\infty\) from AMD2, a rational vector is sent from
the real coordinates \(P_v^{-1}\phi_v(x)\) to its scalar
extension in \(F_\infty\otimes_F E\).  Matrix multiplication,
the Hermitian form, and the span of each initial part of the
frame commute with this map.  The original real \(\texttt{InS}\)
and \(\texttt{InN}\) conditions therefore become the first
and third conditions of (AMD.10).  In the third condition the
extra \(z\)-clause is automatic because \(z\in U^\perp\).
The original \(\texttt{Hset}\), which fixes the union of the
two frame ranges and \(\{z\}\), becomes exactly \(H(F_\infty)\).
Thus \(b_\infty\) restricts to homeomorphic group isomorphisms
on the two specified closed real subgroups.

By AMD6 these subgroups are unimodular, even if \(H(F_\infty)\)
is noncompact.  Apply AMD8 to the two supplied ambient measures
\(\nu_N,\nu_H\), and push their subgroup Haar measures forward
under the two restricted isomorphisms.  Call the results
\(\mu_{N,\infty},\mu_{H,\infty}\).  A topological group
isomorphism carries Haar measure to Haar measure.  By composition
of the two inclusions and the original bridge,
\[
 \operatorname{amb}\mu_{N,\infty}=(b_\infty)_*\nu_N,
 \qquad
 \operatorname{amb}\mu_{H,\infty}=(b_\infty)_*\nu_H.
 \tag{AMD.20}
\]
AMD2 identifies \(b_\infty\) with the source's \texttt{toInfL}
for the supplied witness, not merely a convenient replacement.
Equations (AMD.20), together with the two Haar assertions, are
every field of the original \texttt{InfTransfer}.
\end{proof}
% END AMD9

\subsection{Every field of the original measured block}

% BEGIN AMD10
\paragraph{AMD10. The full global-group measure constructor.}
\label{claim:AMD10}
The original \texttt{MGExists} holds for every \(I\), every
given sigma-finite Haar measure \(\mu_f\) on \(I.G_f\), and
every \(r\in\mathbb N\).  The three measures specified by
its conclusion are the exact pushforwards of the original
global, infinite, and finite measures.

\begin{proof}
Use the one bridge \(b\) from AMD2.  Let \(\mu_\infty\)
be the source's product of the fixed real Haar measures
\(\texttt{REAL.muG}\), so its original full measure is
\(\mu_\infty\times\mu_f\).  Define
\[
 \mu_G=b_*(\mu_\infty\times\mu_f),\qquad
 \mu_{G,\infty}=(b_\infty)_*\mu_\infty,\qquad
 \mu_{G,f}=(b_f)_*\mu_f.
 \tag{AMD.21}
\]
These are Haar measures.  For every finite place choose any
Haar measure \(\mu_{G,v}\) on \(G(F_v)\).  For every tuple
\(X_0\) of length \(2r+2\), choose any Haar measure \(\mu_B(X_0)\)
on its pointwise stabilizer in \(G(\mathbb A_F)\).
AMD4 proves the needed closedness, local compactness, and second
countability even for a degenerate tuple; no nonsingularity
condition on \(X_0\) is imposed here.

These are the five entries of the original \texttt{MeasG}.
All five Haar clauses hold.  The remaining clause follows
from (AMD.5) and product pushforward:
\[
        \operatorname{split}_*\mu_G
            =\mu_{G,\infty}\times\mu_{G,f}.
 \tag{AMD.22}
\]
Thus the entire \texttt{MeasG.IsTextNorm} holds.  Equations
(AMD.21), with \(b_\infty,b_f\) identified in AMD2, are
exactly the three equalities in \texttt{MGExists}.  The freely
chosen local and tuple-stabilizer measures have no further
normalization requirement in this original contract.
\end{proof}
% END AMD10

% BEGIN AMD11
\paragraph{AMD11. The whole generic frame-measure core.}
\label{claim:AMD11}
The full original \texttt{MWCore} holds for every CM extension,
every nondegenerate Hermitian matrix, every original standard
frame, and every pair of given infinity Haar measures
\(\mu_{N,\infty},\mu_{H,\infty}\).  Both infinity measures
are retained as equalities, not merely up to scale.

\begin{proof}
Choose arbitrary Haar measures \(\mu_{N,f}\) and \(\mu_{H,f}\)
on the finite groups.  By the restricted splitting isomorphisms
in AMD4 define the global measures by
\[
\begin{split}
 \mu_N&=(\operatorname{split}_N^{-1})_*
                  (\mu_{N,\infty}\times\mu_{N,f}),\\
 \mu_H&=(\operatorname{split}_H^{-1})_*
                  (\mu_{H,\infty}\times\mu_{H,f}).
\end{split}
 \tag{AMD.23}
\]
They are Haar.  Define \(\mu_{S,f}\) by (AMD.15) from these
two finite measures, and choose arbitrary Haar measures
\(\mu_{S,v}\) on \(S(F_v)\) for every finite place.
These are exactly the eight entries of \texttt{MeasW}.
AMD7 proves the Haar property for \(\mu_{S,f}\); all other
Haar properties follow from the construction.

It remains to check the three actual normalization equalities.
Pushing (AMD.23) into the ambient groups and then splitting gives
\[
\begin{split}
 \operatorname{split}_*(\operatorname{amb}\mu_H)
   &=\operatorname{amb}\mu_{H,\infty}
                       \times\operatorname{amb}\mu_{H,f},\\
 \operatorname{split}_*(\operatorname{amb}\mu_N)
   &=\operatorname{amb}\mu_{N,\infty}
                       \times\operatorname{amb}\mu_{N,f}.
\end{split}
 \tag{AMD.24}
\]
The inclusion squares commute by AMD4, so these are literal
pushforward identities.  The third equality is (AMD.16) over
\(\mathbb A_{F,f}\).  Thus every one of the eleven clauses of
\texttt{MeasW.IsTextNorm} holds.  Its infinity entries were
never changed.  This proves \texttt{MWCore} under precisely
the original standard-frame hypothesis, without asserting the
unguarded auxiliary statement for arbitrary nonstandard frames.
\end{proof}
% END AMD11

% BEGIN AMD12
\paragraph{AMD12. The full original \texttt{MExists} and all same-bridge ties.}
\label{claim:AMD12}
For every original \(D\), the entire \texttt{MExistsCore}\((D)\)
and \texttt{MExists}\((D)\) hold.  For the literal original
antecedent blocks \(A,L,P,R_f,X_u,B_u,\Pi,M_u,W_b\), one
obtains \(M_b\) and a single bridge witness for which the
constructed setting has the complete original \texttt{MeasTie}.

\begin{proof}
First fix \(A,L\), a given finite Haar measure \(\mu_f\),
and the word-indexed weak infinity measures
\(\nu_N(u),\nu_H(u)\) in \texttt{MExistsCore}.  AMD10 gives
one \(\texttt{MeasG}\) and one bridge \(b\) for \(A.I\).
For each word use the original coordinate frame
\[
       P_u=(e_C(A,L),f_C(A,L),z_C(A,L,u)).
 \tag{AMD.25}
\]
This is standard.  Indeed the coordinate map sends the original
Hermitian form to \(A.I.J\); hence the six pairing identities
are exactly the fields \(L.ee,L.ff,L.ef,L.zu\_orth,L.zu\_norm\).
Its scalar is \(T=(D.\mathrm{mem}(u)).A.T\), viewed in \(A.I.F\),
and it is nonzero because the original integer \(T\) is prime.
These verify every field of the source's \texttt{frame\_std}
for this actual frame.

Apply AMD9 with the already chosen bridge and the given weak
measures at this word.  Apply AMD11 to the resulting two
infinity Haar measures, keeping them exactly.  The resulting
\(W_u\) has the complete normalization record and satisfies
\[
 \operatorname{amb}(W_u.\mu_{N,\infty})
                   =(b_\infty)_*\nu_N(u),\qquad
 \operatorname{amb}(W_u.\mu_{H,\infty})
                   =(b_\infty)_*\nu_H(u).
 \tag{AMD.26}
\]
This is also the full original \texttt{MWExistsStd} construction.
Choose such \(W_u\) for every word.  There is no new bridge
choice at a word, since the original bridge was fixed first.
The \(\texttt{MeasG}\) and this family are precisely the
four fields of \(M_b:\texttt{MBlk}\).  Equations (AMD.21)
and (AMD.26) give all three group-measure ties and both
word-indexed subgroup-measure ties in \texttt{MeasTieCore}.
This proves \texttt{MExistsCore} with its full quantifiers.

For \texttt{MExists}, take \(\mu_f=B_u.\mu_f\), with its
original Haar and sigma-finiteness fields, and take the original
\(M_u.\mu_N,M_u.\mu_H\).  Its two weak-Haar fields are
exactly the antecedents just used.  In the source's dependent
constructor \(\texttt{mk}\), the relevant fields are these
same measures and these same coordinate maps.  Consequently
\texttt{MeasTieCore} becomes the literal \texttt{MeasTie}
displayed in \texttt{HidaSetting}, without a change of variables
or a scalar adjustment.  The other supplied blocks, including
\(W_b\), are retained as given.  No gap hypothesis, no Bessel
existence assertion, and no Siegel--Weil formula is needed to
prove this measured-block constructor.
\end{proof}
% END AMD12

\subsection{Prescribing global measures while fixing the infinite factors}

% BEGIN AMD13
\paragraph{AMD13. Absorbing the sole global Haar scalar into the finite factor.}
\label{claim:AMD13}
Let \(e:A\times B\simeq C\) be a topological group
isomorphism of locally compact second-countable groups.  For
every given Haar measure \(\alpha\) on \(A\) and every
given Haar measure \(\gamma\) on \(C\), there is a Haar
measure \(\beta\) on \(B\) such that
\[
                       e_*(\alpha\times\beta)=\gamma.
 \tag{AMD.27}
\]
In particular the original \texttt{exists\_haar\_map\_prod\_eq}
is supplied with the prescribed \(A\)-measure unchanged.

\begin{proof}
Choose any Haar measure \(\beta_0\) on \(B\).  The
pushforward \(\gamma_0=e_*(\alpha\times\beta_0)\) is Haar.
Haar uniqueness gives \(\gamma=c\gamma_0\) for a finite
positive real number \(c\).  Put \(\beta=c\beta_0\).
It is Haar, and product integration and pushforward commute
with this finite scalar multiplication, giving (AMD.27).
The first factor \(\alpha\) is not rescaled.  This argument
also works when one or more of the groups is trivial; a Haar
measure there is a finite positive mass at its sole point.
\end{proof}
% END AMD13

% BEGIN AMD14
\paragraph{AMD14. The entire original prescribed-measure extension.}
\label{claim:AMD14}
For every original standard frame, prescribe Haar measures
\(\mu_{N,\infty},\mu_{H,\infty},\mu_N,\mu_H\) and a
family of Haar measures \(\mu_{S,v}\) on the exact local
groups.  There is an original \(W:\texttt{MeasW}\) with
\(W.\texttt{IsTextNorm}\) which has each of these five
specified entries exactly.  This is the whole original
\texttt{exists\_measW\_prescribed}.

\begin{proof}
Apply AMD13 to \(\operatorname{split}_N^{-1}\) with the
given infinite measure and the given global \(N\)-measure.
This supplies a finite Haar measure \(\mu_{N,f}\) for which
the first identity in (AMD.23) holds with the prescribed
\(\mu_N\).  Apply it separately to
\(\operatorname{split}_H^{-1}\), obtaining \(\mu_{H,f}\)
and the second identity in (AMD.23).  Neither infinite measure
nor either global measure has changed.

Define \(\mu_{S,f}\) by AMD7, and use the supplied family
\(\mu_{S,v}\) as the local entry.  The same proof of the
eleven normalization clauses in AMD11 applies.  In particular
the finite \(S\)-measure is the exact product pushforward of
the two newly chosen finite factors.  The source imposes no
additional equality relating this finite measure to the
independently supplied local \(S\)-measures in this contract.
All five requested equalities therefore hold simultaneously.
\end{proof}
% END AMD14

\subsection{The actual unipotent rational quotient and its covolume}

% BEGIN AMD15
\paragraph{AMD15. A compact set of representatives for the original unipotent quotient.}
\label{claim:AMD15}
For every original standard frame, the rational subgroup
\(N(F)\subset N(\mathbb A_F)\) is exactly the source's
\(\texttt{ratOf}\).  It is countable, discrete and closed,
and there is a compact set \(K\subset N(\mathbb A_F)\)
such that
\[
                    N(\mathbb A_F)=N(F)K.
 \tag{AMD.28}
\]
This supplies the full original \texttt{nsub\_cocompact}.

\begin{proof}
The defining flag equations have coefficients in \(F\), after
writing \(E=F\oplus F\delta\).  A rational matrix whose
adelic extension satisfies these equations already satisfies
them over \(F\), since scalar extension is injective.  The
same applies to its unitary equation.  Thus intersection with
the source's rational ambient unitary subgroup is exactly
\(N(F)\), as asserted.

The polynomial logarithm of AMD6 identifies \(N(F)\) with
a finite-dimensional nilpotent \(F\)-Lie algebra
\(\mathfrak n(F)\), and similarly over \(\mathbb A_F\).
Choose an \(F\)-basis of that Lie algebra.  The diagonal
embedding of \(F\) in \(\mathbb A_F\) is discrete and
closed, so the same holds for \(F^{\dim\mathfrak n}\) and
hence for \(N(F)\).  Countability follows from countability
of the number field.

We give the needed compactness argument explicitly.  In the
additive case choose a closed fundamental parallelepiped for
\(\mathcal O_F\subset F_\infty\) and multiply it by
\(\prod_{v<\infty}\mathcal O_{F,v}\).  This is compact and
its translates by \(F\) cover \(\mathbb A_F\): the Chinese
remainder theorem, after clearing the finitely many denominators,
first makes every finite component integral, and translation by
\(\mathcal O_F\) then moves the infinite component into the
parallelepiped.  This also verifies the asserted discreteness
of \(F\): an infinite neighborhood avoiding the nonzero
integral lattice points, times the integral finite adeles,
meets the diagonal \(F\) only at zero.

Since \(\mathfrak n\) is nilpotent, it has a flag of
\(F\)-ideals
\(0=\mathfrak n_0\subset\mathfrak n_1\subset\cdots
\subset\mathfrak n_d=\mathfrak n\), of successive dimension
one, with \([\mathfrak n,\mathfrak n_i]\subseteq\mathfrak n_{i-1}\).
To construct it, choose a nonzero central \(F\)-vector and
induct in the quotient; its center is nonzero whenever the
nilpotent quotient is nonzero.  Exponentiation gives normal
subgroups \(N_i\), with each quotient the additive group.
Choose \(X_i\in\mathfrak n_i\) mapping to its fixed basis
in \(\mathfrak n_i/\mathfrak n_{i-1}\).  The finite
Baker--Campbell--Hausdorff formula, or induction through these
central quotients, shows that
\[
 (t_1,\ldots,t_i)\longmapsto
             \exp(t_1X_1)\cdots\exp(t_iX_i)
 \tag{AMD.29}
\]
is a polynomial bijection onto \(N_i\), with polynomial
inverse.  All its coefficients are in \(F\), so it applies
both to rational and to adelic points.

Inductively suppose \(N_{i-1}(\mathbb A_F)=N_{i-1}(F)K_{i-1}\)
with \(K_{i-1}\) compact.  Write an element of \(N_i\) as
\(u\exp(tX_i)\), and choose \(a\in F\) with \(t-a\)
in the fixed compact additive representative set \(C\).
Then
\[
 u\exp(tX_i)=\exp(aX_i)
   \bigl(\exp(-aX_i)u\exp(aX_i)\bigr)\exp((t-a)X_i).
 \tag{AMD.30}
\]
The middle factor is in \(N_{i-1}(\mathbb A_F)\), because
that subgroup is normal.  Reduce it by the induction hypothesis.
The resulting left factor is rational and the right factor lies
in the compact set \(K_{i-1}\exp(CX_i)\).  This proves
(AMD.28) at the next stage and hence at \(i=d\).
For the zero-dimensional group the compact set is \(\{1\}\).
\end{proof}
% END AMD15

% BEGIN AMD16
\paragraph{AMD16. The literal covolume-one Haar measure.}
\label{claim:AMD16}
For every original standard frame there is a Haar measure
\(\mu_N\) on \(N(\mathbb A_F)\) with
\(\texttt{covol}(\mu_N)=1\), for exactly the source's
choice-based covolume definition.  This is the whole original
\texttt{exists\_haarMeasure\_}\allowbreak
\texttt{unipotentRadical\_}\allowbreak\texttt{covolume\_one}.

\begin{proof}
Put \(G=N(\mathbb A_F)\), \(\Gamma=N(F)\), and choose
any Haar measure \(\lambda\) on \(G\).  AMD15 shows that
\(\Gamma\backslash G\) is compact.  Choose a relatively
compact open neighborhood \(V\) of the identity with
\(VV^{-1}\cap\Gamma=\{1\}\).  Such a neighborhood exists
by discreteness.  The quotient map is open, so compactness
gives finitely many \(x_1,\ldots,x_m\) such that the images
of \(Vx_j\) cover the quotient.  Set
\[
 B_j=Vx_j\setminus\bigcup_{i<j}\bigcup_{\gamma\in\Gamma}
                                     \gamma Vx_i,
 \qquad B=\bigcup_{j=1}^m B_j.
 \tag{AMD.31}
\]
These are Borel sets, because \(\Gamma\) is countable.
No two points of one \(Vx_j\) have the same orbit, by the
condition on \(VV^{-1}\), and (AMD.31) removes every orbit
already represented at an earlier stage.  Thus \(B\) meets
each orbit exactly once.  It is a Borel fundamental domain
contained in the finite union of the compact sets
\(\overline Vx_j\).  Hence \(\lambda(B)<\infty\).
Also \(\lambda(B)>0\), since the countably many translates
\(\gamma B\) cover \(G\) and a Haar measure is nonzero.

Every measurable fundamental domain has this same measure.
Indeed partition one of them into its intersections with the
countably many translates of the other, translate each piece
back, and use left invariance and countable additivity.  The
same argument works for the source's almost-everywhere
disjointness and covering definition after discarding the
corresponding null sets.  In particular its chosen fundamental
domain has the finite positive measure \(\lambda(B)\).
Multiplying \(\lambda\) by a positive finite constant
preserves the null sets, the fundamental-domain property, and
the preceding equality of measures.

Define \(\mu_N=\lambda(B)^{-1}\lambda\).  It is Haar,
its fundamental domains have measure one, and the real value
of that measure is one.  The source's definition selects the
measure of a fundamental domain and takes its real part when
such a domain exists.  Existence and choice independence have
just been proved, so its value is exactly one for this
\(\mu_N\).  This proves the stated source contract, rather
than only compactness of the rational quotient.
\end{proof}
% END AMD16

% BEGIN AMD17
\paragraph{AMD17. The covolume normalization can coexist with every prescribed infinity and local measure.}
\label{claim:AMD17}
For an original standard frame, fix arbitrary infinity Haar
measures \(\mu_{N,\infty},\mu_{H,\infty}\), an arbitrary
global Haar measure \(\mu_H\), and arbitrary local Haar
measures \(\mu_{S,v}\).  One original normalized
\(W:\texttt{MeasW}\) retains all these entries and has
\(\texttt{covol}(W.\mu_N)=1\).  It can be chosen for each
original word using the same already fixed infinity bridge.

\begin{proof}
Choose the global \(\mu_N\) of AMD16 and apply AMD14 with
this measure and all the specified entries.  Its finite factor
absorbs whatever positive scalar is required; neither infinity
factor, the global \(H\)-measure, nor any local \(S\)-measure
is altered.  The resulting \(W\) has the eleven original
normalization clauses and the required covolume.  If the
infinity measures came from AMD9, their exact ambient ties
(AMD.20) therefore survive this replacement.  Applying the
construction independently at every word changes no bridge
and gives the literal preservation step used in
\texttt{RescaleStep3}.

The other inputs to that original step are distinct: the
compatible global orbital measures and stabilizer measures,
the normalized local fiber measures, and their identities from
the original Siegel--Weil premise.  This claim says that once
those prescribed \(H\)- and local \(S\)-Haar measures have
been produced, the unipotent covolume and all required infinity
ties can be imposed simultaneously.  It does not assume that
the prescribed measures already satisfy any additional orbital
or fiber identity in order to construct \(W\).
\end{proof}
% END AMD17

\subsection{Exact effect on the remaining global construction}

AMD12 removes the entire measured-block producer from the
remaining work, and AMD3 supplies the full adelic isomorphism,
including the rational and compact subsets.  Together with the
eight dependent blocks in ASG18, this gives every original
setting block except the normalized nonzero Bessel block
\(W_b\).  For any such original \(W_b\), the dependent
constructor now gives an actual \(X\) with its complete
\texttt{MeasTie}.  The decay assertion associated with a
properly constructed \(W_b\) is not inferred from measure
existence.

The later \texttt{RescaleStep3} still needs its original
Siegel--Weil and orbital/fiber comparison suppliers.  AMD14,
AMD16, and AMD17 reconstruct its prescribed-measure and
unipotent-normalization portions while preserving the original
infinity measures.  The Witt/genus/level/representation
transport into the arithmetic columns also remains a separate
same-object construction.  None of these distinctions closes
\texttt{GLOBAL\_SPECTRAL}, the arithmetic counting obligation,
or the full RH/GRH reconstruction by itself.
